% =============================================================================
% Chapter 8: Verification, Milestones, and Possible Outcomes
% Input-ready fragment for the lean toy-model research proposal.
% This file intentionally contains no preamble or document environment.
% =============================================================================

\chapter{Verification, Milestones, and Possible Outcomes}
\label{ch:verification-milestones-outcomes}

The preceding chapters define the proposed ontology, the sector calculations,
the common-medium test, and the later throat and reservoir program. This final
chapter states what evidence permits any of those claims to advance and what
conclusion is authorized at each stage. It does not reproduce the detailed
artifact schemas, discrepancy taxonomy, numerical procedures, or revision
records maintained in the derivation ledger. Its purpose is narrower: to keep
the top-down proposal, the bottom-up calculations, and the scientific decision
sequence attached to the same controlled candidate.

A result belongs to the current program only when the question, assumptions,
branch, frozen inputs, calculation, evidence, and downstream use all refer to
the same version. A recent file, a successful script run, or a plausible plot is
not enough. Conversely, a negative or inconclusive result is not discarded
merely because it interrupts the intended sequence. The verification process is
part of the falsification test because it determines whether several apparent
sector successes truly arise from one medium.

\section{Proposal-to-ledger verification}

The proposal and the derivation ledger are independent views of the same
research program. The proposal works from observable behavior downward: it
states what must be reproduced, which unwanted channels must be included, what
the result would require from the common medium, and what would reject the
branch. The derivation ledger works from the candidate equations upward: it
records the selected branch and boundary data, performs the symbolic and
numerical calculations, preserves the actual signs and extra modes, and tracks
the inputs and outputs inherited by later steps.

These views should be compared through a concise proposal-to-ledger crosswalk.
For every mandatory requirement, the crosswalk should identify the exact claim,
the Candidate Parent System version and dynamical branch, the background and
environment, the boundary or interaction ensemble, all frozen inputs, the valid
scale and approximation regime, the ledger steps that test it, the human and
executable evidence, the expected and observed result, the current decision,
and the downstream dependencies. One requirement may need several calculations,
and one calculation may support several requirements; the mapping therefore
must work in both directions.

The minimum completeness rule is straightforward. No mandatory proposal claim
should lack either a current evidence path or an explicit open-status record.
No promoted ledger result should lack a stated proposal purpose and claim
boundary. No derived output should be used downstream without a dependency
record, and no current pass should rely on a superseded candidate, an unstated
ensemble, or a reopened upstream result.

This comparison is also where discrepancies become scientifically useful. The
proposal is not rewritten automatically to match an unexpected output, and a
complicated expression is not declared successful merely because the ledger
produced it. A mismatch may reveal that the proposal asked the wrong physical
question, that the derivation used different assumptions, that a convention or
implementation is inconsistent, or that the candidate fails. Until the earliest
cause is identified and the affected evidence is rerun, the claim remains open
or reopened.

Current calculation status is therefore read from the derivation ledger,
executable artifacts, and the independent Codex repository audit rather than from a
dated status chapter in this proposal. The proposal records the enduring
research contract; the audit records which parts of that contract presently
have current evidence.

\section{Required symbolic and numerical evidence}

Every load-bearing claim begins with a human-readable derivation. It should state
the governing equations and whether each is postulated or inherited, the field
variables and units, the conservative, relaxational, or mixed branch, all
constraints and symmetries, the boundary or reservoir data, the approximation
order, and the conclusion the calculation is authorized to support. Intermediate
steps must remain visible where signs, cancellations, singular limits, source
normalizations, or branch choices affect the result. A necessary-condition
calculation obtained by solving for a desired coefficient may be useful, but it
must not be reported as a derivation of that coefficient from the candidate.

Where symbolic verification is feasible, SymPy and Mathematica should construct
the central mathematics independently from the same controlled specification.
They may share the symbol dictionary, assumptions, and expected invariants, but
one implementation should not import the other's final symbolic expression.
Comparison should use canonical objects rather than visual similarity: equation
residuals, characteristic polynomials, root multiplicities, invariant
subspaces, projectors, ranks, principal minors, series coefficients, symmetry
transforms, current divergences, and high-precision samples at regular and
near-threshold points. Exact arithmetic and explicit assumptions should be
retained as long as practical.

Agreement between the two systems supports the algebraic claim under the stated
assumptions. It does not establish that the parent postulates are true, that the
required nonlinear branch exists, or that the result is observationally
acceptable. If one system cannot complete a calculation, the missing check may
be replaced by an independent hand reduction, another mathematical method, a
controlled high-precision residual, or an alternative numerical implementation,
but the asymmetry must remain visible.

Numerical evidence is required for most existence, localization, stability, and
closure claims. A stable finite slab, guided or resonant modes, a localized
photon candidate, a supported throat, moving and accelerating families,
two-core contact, and the local--global fixed point cannot be established by a
formal symbolic expression alone. The numerical record must retain the exact
equations and baseline, domain and boundary conditions, discretization and
solver, initial data or continuation seed, raw solution fields, residuals and
constraint violations, and the observables used by the gate.

A gate-level numerical result also requires the refinements relevant to its
physical classification. These may include grid or basis refinement, timestep
and domain variation, outgoing-boundary or regulator variation, precision
increase, continuation in both directions, multiple seeds, and comparison of
independent solvers or discretizations at selected points. Conservation,
supplied and dissipated power, radiation into every coupled branch, entropy and
reservoir balance, and linear, Floquet, or retarded-response stability must be
checked when the chosen branch requires them. Open, dissipative, or non-normal
problems may additionally require complex poles, resonance widths, residues,
passivity, and sensitivity to transient growth.

The evidence standard must match the mathematical problem. A static energy
minimum cannot validate a maintained dissipative state. A conservative Floquet
spectrum cannot close a mixed branch whose reservoir was omitted. A reduced
envelope or imposed throat geometry is a precursor until it is continued to the
parent fields. Failure of one solver or one seed is not a proof of nonexistence;
one converged solution is not proof of uniqueness, robustness, or an acceptable
family.

\section{Milestones and decision gates}

The gates follow the mathematical closure order, even when conditional work is
scheduled in parallel. Each exit condition authorizes the next level of claim;
it does not confer success on later stages. The lean program uses the sequence
in \cref{tab:program-gates}.

\begingroup
\small
\begin{longtable}{@{}L{0.08\textwidth}L{0.27\textwidth}L{0.57\textwidth}@{}}
\caption{Major verification and decision gates.}
\label{tab:program-gates}\\
\toprule
\textbf{Gate} & \textbf{Milestone} & \textbf{Minimum exit condition} \\
\midrule
\endfirsthead
\toprule
\textbf{Gate} & \textbf{Milestone} & \textbf{Minimum exit condition} \\
\midrule
\endhead
G0
& Controlled research contract
& Canonical definitions, proposal requirements, failure conditions, shared
  quantity identities, and the proposal--ledger crosswalk are current and
  mutually consistent. \\
\addlinespace
G1
& Light and far-field capability contracts
& The light, gravity, electric, magnetic, radiation, mass-response, extra-mode,
  leakage, and drag programs have produced explicit common-medium requirements
  and named failure channels rather than separate tuned solutions. \\
\addlinespace
G2
& Common region and Candidate V1
& The Medium Capability Matrix identifies a nonempty contract-feasible region,
  and one Candidate Parent System version freezes its fields, branch, equations,
  operator basis, coefficients, symmetries, source conventions, and boundary and
  reservoir data. \\
\addlinespace
G3
& Stable background and common far field
& Candidate V1 produces a selected stable or sufficiently metastable finite
  brane, a complete acceptable constrained spectrum, and all mandatory declared
  light and far-field capabilities at the same current point or region. \\
\addlinespace
G4
& Supported throat and species basis
& At least one self-consistent stable or sufficiently long-lived supported
  oriented throat exists; its source profiles replace the provisional sources;
  and the required orientation-partner, species, charge, and conservation tests
  are defined and passed to the declared scope. \\
\addlinespace
G5
& Moving and accelerating throat closure
& The solved throat family yields positive ordinary passive and inertial
  response, acceptable low-energy universality, one electric--magnetic--radiative
  normalization, and bounded extra radiation, leakage, wake, and drag. \\
\addlinespace
G6
& Two-throat and contact closure
& The same solved objects recover the far field, survive the finite-thickness
  crossover, and produce conservation-respecting contact, annihilation,
  exclusion, or neutral-binding outcomes with same-species partners distinguished
  from different species. \\
\addlinespace
G7
& Localized throat drainage, distributed return, and reservoir closure
& The local one-way throat problem and the global bulk-reservoir problem,
  including distributed porous-brane return through the positive bulk-to-brane
  contribution to \(S_{\mathrm{leak}}\), converge to a stable or sufficiently
  long-lived common fixed point with material, momentum, energy, and, where
  applicable, entropy and internal-reservoir accounting closed. \\
\bottomrule
\end{longtable}
\endgroup

The localized photon-packet track is a stronger light milestone that branches
from the stable-background and complete-spectrum result. It may be pursued in
parallel with later realization work, but it receives its own controlled
existence, stability, leakage, and collision decision. Failure of that stronger
track narrows the outcome to guided classical light unless the frozen candidate
explicitly made a self-bound packet mandatory; it does not silently convert an
unsolved packet into a passed light result.

A gate may be recorded as \emph{pass}, \emph{conditional pass}, \emph{reopen},
\emph{fail}, or \emph{open/inconclusive}. A conditional pass must name the
unresolved condition and the limited reason downstream exploratory work remains
useful; it cannot conceal a known violation of a mandatory requirement. A fail
means the current branch or candidate misses a frozen condition. An open result
means that evidence, convergence, domain coverage, or the mathematical object
itself is still insufficient to decide.

\section{Revision and downstream reopening}

A substantive upstream change is a scientific change, not a prose correction.
If a field, constitutive term, coefficient, dynamical branch, source definition,
calibration, boundary ensemble, or reservoir law changes, the candidate receives
a new version and every actual dependent result is marked reopened until it is
rederived. Derived outputs are frozen downstream inputs: a guided-light profile,
electric eigenvector, throat residue, support frequency, or reservoir state
cannot become adjustable again merely because a later sector prefers another
value.

Reopening should follow real dependencies rather than trigger ceremonial reruns
of unrelated work. A revised shear law or slab-thickness mechanism reopens the
background spectrum, guided light, electric response, characteristic speeds,
and throat support. A revised oriented-current normalization reopens static
charge, magnetism, transverse radiation, and charge universality. A revised
order or reservoir branch reopens gravity, throughput, throat lifetime, inertia,
drag, and global closure. A revised antiparticle map or physical core diagnostic
reopens species, pair, contact, and annihilation claims without necessarily
invalidating the homogeneous background spectrum.

Historical calculations remain valuable as failed branches, special limits,
regression tests, or sources of new requirements, but they no longer count as
current evidence for the revised candidate. A new mechanism introduced after a
failure is permitted only as an explicit revision of the common parent system,
with its consequences propagated through every affected sector. It may not be
added privately to the calculation that exposed the problem.

\section{Interpreting success, partial survival, and failure}

The outcome of the program is layered. Every report should state the strongest
coherent submodel that survives, the exact gate reached, the gate that failed or
remains open, and the assumptions under which the conclusion holds.

\begin{center}
\begin{tabularx}{\textwidth}{@{}L{0.27\textwidth}X@{}}
\toprule
\textbf{Outcome} & \textbf{Authorized conclusion} \\
\midrule
Open or inconclusive
& A required equation, branch, source, convergence study, acceptance bound, or
  domain test is missing or conflicting. Neither pass nor failure is justified. \\
Sector result
& One named behavior, compatibility condition, or no-go statement is established
  within its assumptions. It need not establish any larger part of the ontology. \\
Partial analog
& A coherent subset survives after a stronger claim fails or remains outside
  scope, such as guided classical light without a self-bound photon packet or a
  local gravity analog with a limited return regime. \\
Common far-field medium
& One frozen candidate has a stable background and rederives all mandatory
  declared light and far-field sectors at one shared point or region. Particle
  existence and global closure remain unproved. \\
Supported particle sector
& One or more stable or sufficiently long-lived throat branches exist with
  solved source profiles and the declared orientation, species, and charge
  structure. Moving, pair, or reservoir gates may still fail. \\
Moving and interacting particle sector
& The candidate also closes ordinary response, universality, magnetism,
  radiation, drag, two-body crossover, contact, and the declared short-range
  outcomes. The global reservoir fixed point may remain open. \\
Full declared classical survival
& One current parent system passes every mandatory gate through the local--global
  fixed point in the stated regimes, with all shared inputs frozen and all
  unwanted channels retained in the evidence. \\
Candidate failure
& One named Candidate Parent System violates a mandatory gate or cannot close
  its own background, throat, interaction, or reservoir problem. Other admitted
  candidates are not thereby rejected. \\
Architecture failure
& The mandatory intersection is shown to be empty throughout the admitted
  architecture, an essential mechanism is impossible across its permitted
  realizations, or success requires prohibited sector-specific retuning or an
  unsupported new mechanism. \\
\bottomrule
\end{tabularx}
\end{center}

Meaningful negative results include several important intermediate outcomes.
Classical guided light may survive while every permitted self-binding mechanism
for a free photon packet fails. Gravity and electricity may work in separate
effective limits but have no common coefficient region. A stable common brane
and far field may exist while no self-consistent supported throat does. Stable
throats may exist while unavoidable drag, excess radiation, wrong-sign response,
or nonuniversal charge or mass ratios reject the particle interpretation. Local
particle physics may close while the global distributed-return and reservoir
iteration has no conservation-respecting fixed point. Each result narrows the
viable model and identifies where the one-medium program stops.

Full declared survival would be the strongest result available to this proposal,
but its scope remains classical and conditional on the stated regimes. It would
not establish that the real universe is literally this medium, reproduce exact
general relativity or exact Maxwell theory, or supply quantum statistics, a
universal quantum of action, or a complete particle spectrum. A later
cosmological extension would remain a separate result unless deliberately made
mandatory. The value of the program is therefore not limited to a final positive
answer: a controlled partial analog, a precise candidate failure, an empty
common region, or a well-supported no-go result is also a successful scientific
outcome because it is obtained without hiding dependencies or repairing the
model after the fact.
